\documentclass[a4paper,12pt]{article}
\usepackage[T2A]{fontenc}
\usepackage[utf8]{inputenc}
\usepackage[russian]{babel}
\usepackage{amsmath,amssymb,icomma}
\usepackage[a4paper,left=20mm,right=20mm,top=21mm,bottom=22mm,headheight=25pt]{geometry}
\usepackage{microtype,booktabs,tabularx,array,enumitem,xcolor,tikz,fancyhdr,setspace,parskip,titlesec}
\usetikzlibrary{arrows.meta,positioning,shapes.geometric}
\usepackage[hidelinks]{hyperref}
\definecolor{accent}{HTML}{008080}
\definecolor{darkaccent}{HTML}{004D4D}
\definecolor{textgray}{HTML}{333333}
\color{textgray}
\onehalfspacing
\setlength{\parskip}{5pt}
\setlength{\parindent}{0pt}
\setlist[itemize]{leftmargin=*,itemsep=1pt,topsep=3pt}
\setlist[enumerate]{leftmargin=*,itemsep=8pt,topsep=5pt}
\titleformat{\section}{\Large\color{darkaccent}}{\thesection.}{0.5em}{}[\vspace{4pt}\color{accent}\titlerule]
\titlespacing*{\section}{0pt}{4pt}{12pt}
\titleformat{\subsection}{\large\color{darkaccent}}{}{0pt}{}
\titlespacing*{\subsection}{0pt}{9pt}{3pt}
\pagestyle{fancy}
\fancyhf{}
\fancyhead[L]{\small Дроби, пропорции, проценты и масштаб}
\fancyhead[R]{\small Ксения Васильченко}
\fancyfoot[C]{\small\thepage}
\renewcommand{\headrulewidth}{0.5pt}
\renewcommand{\headrule}{\hbox to\headwidth{\color{accent}\leaders\hrule height \headrulewidth\hfill}}
\newcommand{\idea}[1]{\textcolor{darkaccent}{\textbf{#1}}}
\newcommand{\ex}[1]{\par\medskip\textcolor{darkaccent}{#1}\par\smallskip}
\newcommand{\cm}{\,\text{см}}
\newcommand{\km}{\,\text{км}}
\newcommand{\kg}{\,\text{кг}}
\newcommand{\mtr}{\,\text{м}}
\newcommand{\pct}{\,\%}
\begin{document}
\begin{titlepage}
\centering
\vspace*{28mm}
{\Large\color{darkaccent} МАТЕМАТИКА\par}
\vspace{16mm}
{\Huge Чек-лист\par}
\vspace{9mm}
{\LARGE\color{accent} Десятичные дроби, уравнения,\\[2pt] пропорции, проценты и масштаб\par}
\vspace{15mm}
{\large Подготовка к обзорной контрольной работе\par}
\vfill
{\large Ксения Васильченко\par}
\vspace{9mm}
{26 ноября 2026 г.\par}
\vspace{15mm}
{\small Лёвин Артём Александрович\\эксклюзивно для Ксении Васильченко\par}
\vspace{18mm}
\begin{tikzpicture}
\draw[accent,line width=.9pt] (0,0) .. controls (2.1,0.8) and (3.9,-0.8) .. (6,0) .. controls (8.1,0.8) and (9.9,-0.8) .. (12,0);
\fill[accent] (6,0) circle (1.1pt);
\end{tikzpicture}
\vspace*{3mm}
\end{titlepage}
\setcounter{page}{1}
\section{Десятичные дроби: как не потерять запятую}
\idea{Главная идея.} При вычислениях важно сохранять разряды числа: единицы, десятые, сотые, тысячные.
\subsection{Умножение и деление на 10, 100, 1000 и т. д.}
При умножении на единицу с нулями запятая перемещается вправо, при делении --- влево. Число шагов равно числу нулей в множителе или делителе. Если цифр не хватает, допишите нули.
\begin{center}
\begin{tikzpicture}[>=Stealth, every node/.style={font=\normalsize}]
\node (c) at (0,0) {$1{,}25$};
\node (l) at (-4.2,0) {$0{,}00125$};
\node (r) at (4.2,0) {$1250$};
\draw[->,accent,thick] (c.west) -- node[above] {$:1000$} (l.east);
\draw[->,accent,thick] (c.east) -- node[above] {$\cdot1000$} (r.west);
\end{tikzpicture}
\end{center}
\[
0{,}0285\cdot100=2{,}85,\qquad 0{,}1234\cdot1000=123{,}4.
\]
\subsection{Деление на десятичную дробь}
Чтобы делитель стал целым числом, умножьте и делимое, и делитель на одно и то же число (10, 100, 1000 и т. д.). Частное не изменится:
\[
160{,}82:4{,}3=\frac{160{,}82\cdot10}{4{,}3\cdot10}
=\frac{1608{,}2}{43}=37{,}4.
\]
При делении столбиком поставьте запятую в частном, когда переходите к цифрам дробной части делимого. Проверка: $37{,}4\cdot4{,}3=160{,}82$.
\subsection{Сложение и вычитание}
Записывайте запятую под запятой, одинаковые разряды --- друг под другом. При необходимости дописывайте нули справа: $0{,}57+0{,}86=1{,}43$.
\ex{Пример. Вычитание столбиком}
\[
\begin{array}{r@{}l}
 &37{,}4\\[-2pt]
-&25{,}3\\\hline
 &12{,}1
\end{array}
\qquad 12{,}1+25{,}3=37{,}4.
\]
\idea{Самопроверка.} При умножении положительного числа на 100 оно увеличивается в 100 раз, а при делении --- уменьшается. Запятая перемещается ровно на два разряда, не на три.
\clearpage
\section{Уравнения: неизвестный компонент действия}
Рассматривайте выражение в скобках как единое число, пока не найдёте его значение. Начинайте с последнего действия в левой части.
\begin{center}
\renewcommand{\arraystretch}{1.35}
\begin{tabularx}{\linewidth}{lXX}
\toprule
Исходная запись & Найти $a$ & Найти $b$\\\midrule
$a+b=c$ & $a=c-b$ & $b=c-a$\\
$a-b=c$ & $a=c+b$ & $b=a-c$\\
$a\cdot b=c$ & $a=c:b$, если $b\ne0$ & $b=c:a$, если $a\ne0$\\
$a:b=c$ & $a=c\cdot b$ & $b=a:c$, если $c\ne0$\\\bottomrule
\end{tabularx}
\end{center}
В последней строке обязательно $b\ne0$. В случае нулевого множителя или нулевого частного нельзя делить на нуль: такие ситуации проверяют отдельно.
\ex{Пример 1. Неизвестное слагаемое в скобках}
Решите уравнение $(z+25{,}3)\cdot4{,}3=160{,}82$.
Сначала найдите значение всей скобки как неизвестного множителя:
\[
z+25{,}3=160{,}82:4{,}3=37{,}4.
\]
Теперь найдите неизвестное слагаемое:
\[
z=37{,}4-25{,}3=12{,}1.
\]
Проверка: $(12{,}1+25{,}3)\cdot4{,}3=160{,}82$.
\ex{Пример 2. Неизвестный делитель и вычитаемое}
Если $18{,}6:t=3$, то $t=18{,}6:3=6{,}2$.
Если $9{,}4-u=2{,}7$, то $u=9{,}4-2{,}7=6{,}7$.
\ex{Пример 3. Сначала вычислите числовую часть}
\[
y=\frac{0{,}0285\cdot100}{5}+0{,}86
=\frac{2{,}85}{5}+0{,}86=1{,}43.
\]
\idea{Проверяйте.} Подставьте найденное число в исходное уравнение. Левая и правая части должны совпасть. Не меняйте числа и знаки при переписывании.
\clearpage
\section{Деление с остатком и смешанные числа}
Чтобы выделить целую часть неправильной дроби, разделите числитель на знаменатель с остатком. Частное становится целой частью, остаток --- числителем дробной части, знаменатель остаётся прежним.
\[
a=bq+r,\quad 0\le r<b,\quad b>0
\qquad\Longrightarrow\qquad
\frac ab=q+\frac rb.
\]
Здесь $a$ --- неотрицательное целое число, $b$ --- натуральное число.
\ex{Пример. Запишите $31:9$ в виде смешанного числа.}
\[
31=9\cdot3+4\quad\Longrightarrow\quad \frac{31}{9}=3\frac49.
\]
Если остаток равен нулю, пишите просто целое число: $81:9=9$. Если дробная часть сокращается, сократите её: $42:8=5\frac28=5\frac14$.
\section{Прямая пропорциональность}
Две положительные величины прямо пропорциональны, если при увеличении одной в несколько раз другая увеличивается во столько же раз. Их отношение постоянно.
\[
\frac ab=\frac cd\quad\Longleftrightarrow\quad ad=bc,\qquad b\ne0,\ d\ne0.
\]
При составлении таблицы разместите одноимённые величины в одном столбце, а соответствующие друг другу значения --- в одной строке.
\begin{center}
\begin{tikzpicture}[>=Stealth]
\node at (-1.9,1.3) {Величина 1};\node at (1.9,1.3) {Величина 2};
\node (a) at (-1.9,.45) {$a$};\node (b) at (1.9,.45) {$b$};
\node (c) at (-1.9,-.65) {$c$};\node (d) at (1.9,-.65) {$d$};
\draw[accent,thick] (a.east)--(d.west);\draw[accent,thick] (b.west)--(c.east);
\draw[->,accent] (-3,.55)--(-3,-.75);\draw[->,accent] (3,.55)--(3,-.75);
\node at (0,-1.45) {$ad=bc$};
\end{tikzpicture}
\end{center}
\idea{Главная проверка.} Стрелки в одну сторону и перемножение крест-накрест применимы здесь именно потому, что зависимость прямая.
\clearpage
\section{Проценты: целое, часть, процент}
Один процент --- одна сотая часть числа: $1\pct=\frac1{100}$. Всё количество принимается за $100\pct$. Сначала определите, что является целым, а что --- его частью.
\begin{center}
\renewcommand{\arraystretch}{1.3}
\begin{tabular}{lcc}\toprule & Количество & Проценты\\\midrule Целое & $A$ & $100$\\ Часть & $B$ & $p$\\\bottomrule\end{tabular}
\end{center}
При $A>0$ и $0\le B\le A$ из пропорции получаем $100B=Ap$.
\[
B=\frac{Ap}{100},\qquad A=\frac{100B}{p}\ (p>0),\qquad p=\frac{100B}{A}.
\]
\ex{Пример 1. Найти часть по проценту}
В 3,5 кг золотого сплава содержится 41,5\% примесей. Найдите массу примесей, округлив до десятых долей килограмма.
\[
\frac{3{,}5}{x}=\frac{100}{41{,}5}\quad\Longrightarrow\quad
100x=3{,}5\cdot41{,}5=145{,}25.
\]
\[
x=1{,}4525\kg\approx1{,}5\kg.
\]
\ex{Пример 2. Найти целое по его части}
Оливки содержат 64,8\% масла по массе. Сколько килограммов оливок содержат 40,5 кг масла?
\[
\frac{x}{40{,}5}=\frac{100}{64{,}8}\quad\Longrightarrow\quad
x=\frac{40{,}5\cdot100}{64{,}8}=62{,}5\kg.
\]
\ex{Пример 3. Найти процентное содержание}
В пачке творога массой 250 г содержится 41,7 г белка. Какую долю в процентах составляет белок?
\[
\frac{250}{41{,}7}=\frac{100}{x}\quad\Longrightarrow\quad
250x=4170,\quad x=16{,}68\pct.
\]
\idea{Проверяйте.} Масса части не превышает массу целого, а содержание компонента находится от $0\pct$ до $100\pct$. Округляйте только конечный результат, если этого требует условие.
\clearpage
\section{Масштаб: карта и местность}
Масштаб карты $1:n$ показывает, что одной единице длины на карте соответствуют $n$ таких же единиц на местности. Для обычной уменьшенной карты $n>1$.
\[
1\mtr=100\cm,\qquad 1\km=1000\mtr=100\,000\cm.
\]
\begin{center}
\renewcommand{\arraystretch}{1.2}
\begin{tabular}{cc}\toprule Карта (см) & Местность (см)\\\midrule $1$ & $n$\\ $l$ & $L$\\\bottomrule\end{tabular}
\end{center}
При положительных длинах $L=ln$. Отсюда: $l=L:n$, $n=L:l$.
\ex{Пример 1. Из карты в действительность}
Отрезок на карте масштаба $1:10\,000\,000$ равен 4 см. Найдите соответствующее расстояние на местности.
\[
L=4\cdot10\,000\,000=40\,000\,000\cm=400\km.
\]
\ex{Пример 2. Из действительности на карту}
Расстояние на местности равно 5,2 км. Какой длины будет отрезок на карте масштаба $1:1\,000\,000$?
\[
5{,}2\km=520\,000\cm,\qquad
l=\frac{520\,000}{1\,000\,000}=0{,}52\cm.
\]
\ex{Пример 3. Найти сам масштаб}
На местности расстояние от рябины до тайника равно 3,5 м, а на плане --- 4 см. Определите масштаб плана.
\[
3{,}5\mtr=350\cm,\qquad 4n=350,\qquad n=87{,}5.
\]
Ответ: масштаб $1:87{,}5$ (то есть 1 см на плане соответствует 87,5 см на местности).
\idea{Не забывайте единицы.} Если в таблицу вписать километры как сантиметры, ответ будет ошибочным в $100\,000$ раз. Сначала переводите длины в одинаковые единицы, а в конце --- в единицы, которые просит условие.
\clearpage
\section*{Алгоритм: проценты или масштаб}
\thispagestyle{empty}
\vspace{8mm}
\begin{center}
\begin{tikzpicture}[node distance=9mm and 12mm, >=Stealth,
actionnode/.style={draw=accent,rounded corners=3pt,align=center,text width=10.4cm,minimum height=13mm,inner sep=5pt},
term/.style={ellipse,draw=accent,align=center,text width=4.5cm,minimum height=10mm},
branch/.style={diamond,draw=accent,aspect=2.6,align=center,inner sep=2pt,text width=4.1cm},
arrow/.style={->,draw=darkaccent,line width=.8pt}]
\node[term] (start) at (0,0) {Прочитайте условие};
\node[actionnode] (identify) at (0,-1.9) {Определите величины и неизвестное.\\Проверьте, что зависимость прямая.};
\node[branch] (choice) at (0,-4.0) {Задача на\\проценты?};
\node[actionnode,text width=6.1cm] (percent) at (-3.75,-7.0) {Целое --- $100\pct$,\\часть --- $p\pct$.\\Единицы количества одинаковы.};
\node[actionnode,text width=6.1cm] (scale) at (3.75,-7.0) {Масштаб $1:n$: карта --- 1,\\местность --- $n$.\\Все длины в одних единицах.};
\node[actionnode] (table) at (0,-9.7) {Заполните таблицу: три известных числа и $x$.\\Составьте равенство отношений.};
\node[actionnode] (cross) at (0,-11.8) {Перемножьте крест-накрест.\\Решите полученное уравнение.};
\node[actionnode] (check) at (0,-13.9) {Проверьте единицы, смысл и порядок величины.\\При необходимости округлите ответ.};
\node[term] (end) at (0,-15.8) {Запишите ответ};
\draw[arrow] (start)--(identify);\draw[arrow] (identify)--(choice);
\draw[arrow] (choice.west) -| node[pos=.25,above]{Да} (percent.north);
\draw[arrow] (choice.east) -| node[pos=.25,above]{Нет} (scale.north);
\draw[arrow] (percent.south) -- ++(0,-.55) -| ([xshift=-3.75cm]table.north);
\draw[arrow] (scale.south) -- ++(0,-.55) -| ([xshift=3.75cm]table.north);
\draw[arrow] (table)--(cross);\draw[arrow] (cross)--(check);\draw[arrow] (check)--(end);
\end{tikzpicture}
\end{center}
\clearpage
\section*{Тренировочный блок}
\textcolor{darkaccent}{10 заданий от базовых к более сложным.}
Записывайте вычисления последовательно. В задачах на проценты и масштаб составьте таблицу величин и пропорцию, укажите единицы и проверьте смысл результата.
\subsection{Базовый уровень}
\begin{enumerate}
\item При подготовке к контрольной работе нужно восстановить правило перемещения запятой. Вычислите значения четырёх выражений без калькулятора:
\[
0{,}0367\cdot100;\quad 2{,}48:1000;\quad
0{,}1725\cdot1000;\quad 45{,}6:100.
\]
\item Для изготовления одинаковых поделок имеются 47 м ленты. Если разделить её поровну между 8 участниками кружка, сколько метров ленты получит каждый? Запишите результат в виде смешанного числа, выделив целую часть делением с остатком.
\item В школьном задании требуется найти неизвестный компонент действия. Решите уравнение $(x+14{,}6)\cdot2{,}5=53{,}5$. Сначала найдите значение выражения в скобках, затем неизвестное число. Проверьте ответ подстановкой.
\item Для приготовления напитка взяли 18,9 л воды и разлили её поровну в ёмкости. В каждую ёмкость налили по 2,7 л. Обозначьте число ёмкостей буквой $t$, составьте и решите уравнение с неизвестным делителем.
\end{enumerate}
\subsection{Средний уровень}
\begin{enumerate}[start=5]
\item В упаковке ореховой смеси массой 2,8 кг доля миндаля составляет 37,5\% от всей массы. Сколько килограммов миндаля содержится в упаковке? Округлите ответ до десятых долей килограмма.
\item На туристической карте масштаба $1:250\,000$ расстояние между двумя населёнными пунктами изображено отрезком длиной 6,4 см. Найдите соответствующее расстояние на местности. Ответ выразите в километрах.
\end{enumerate}
\clearpage
\section*{Тренировочный блок (продолжение)}
\subsection{Уровень выше среднего}
\begin{enumerate}[start=7]
\item На предприятии перерабатывают семена подсолнечника. Масса масла, содержащегося в семенах, составляет 42\% от массы семян. Какую массу семян необходимо взять, чтобы в них содержалось 31,5 кг масла? Ответ выразите в килограммах. Потери при переработке не учитывайте.
\item На этикетке упаковки творога указано, что масса творога без упаковки составляет 240 г, а масса содержащегося в ней белка --- 39,6 г. Определите процентное содержание белка в твороге. Результат запишите в процентах.
\item Длина участка железной дороги на местности составляет 432 км. Для школьного проекта этот участок нужно изобразить на двух картах: первая имеет масштаб $1:20\,000\,000$, вторая --- $1:25\,000\,000$. Какой длины в сантиметрах будет соответствующий отрезок на каждой карте? Сравните результаты и укажите, на какой карте отрезок длиннее.
\end{enumerate}
\subsection{Повышенный уровень}
\begin{enumerate}[start=10]
\item Ребята составили план школьного двора. Расстояние от входа в школу до дерева на местности равно 7,5 м, а на плане оно изображено отрезком длиной 6 см. Определите масштаб плана в виде $1:n$. Затем найдите, каким отрезком на этом же плане будет изображена дорожка длиной 12,5 м. Длину отрезка выразите в сантиметрах.
\end{enumerate}
\vspace{10mm}
\idea{Перед сдачей работы проверьте:}
\begin{itemize}[label=$\square$]
\item Запятая при умножении и делении перемещена в правильном направлении.
\item В уравнениях найден именно нужный компонент, выполнена подстановка.
\item В задачах на проценты за 100\% принято целое.
\item В задачах на масштаб длины приведены к одинаковым единицам.
\item Записаны ответы с единицами и выполнено требуемое округление.
\end{itemize}
\clearpage
\section*{Ответы}
\begin{center}
\renewcommand{\arraystretch}{1.32}
\begin{tabularx}{\linewidth}{lX}
\toprule № & Ответ\\\midrule
1 & $3{,}67$; $0{,}00248$; $172{,}5$; $0{,}456$\\
2 & $5\frac78\mtr$\\
3 & $x=6{,}8$\\
4 & $18{,}9:t=2{,}7$; $t=7$ ёмкостей\\
5 & $1{,}1\kg$\\
6 & $16\km$\\
7 & $75\kg$\\
8 & $16{,}5\pct$\\
9 & $2{,}16\cm$ и $1{,}728\cm$; на первой карте длиннее\\
10 & $1:125$; $10\cm$\\\bottomrule
\end{tabularx}
\end{center}
\vspace{6mm}
\begin{center}{\large\color{darkaccent} Итоговая памятка}\end{center}
\begin{center}
\renewcommand{\arraystretch}{1.35}
\begin{tabularx}{\linewidth}{p{.31\linewidth}X}\toprule
Сигнал в задаче & Полезное действие\\\midrule
Умножение или деление на $10^k$ & Переместите запятую на $k$ разрядов вправо или влево соответственно.\\
Десятичный делитель & Умножьте делимое и делитель на одно и то же число, чтобы делитель стал целым.\\
Неизвестное в скобках & Найдите значение всей скобки, затем неизвестное внутри неё.\\
Неправильная дробь & Разделите числитель на знаменатель с остатком.\\
Процентное содержание & Определите целое ($100\pct$), часть и неизвестное.\\
Масштаб $1:n$ & Приведите длины к одной единице измерения; составьте пропорцию.\\\bottomrule
\end{tabularx}
\end{center}
\end{document}
